\documentclass[11pt,letterpaper]{article}

\usepackage[top=1in, bottom=1in, left=1in, right=1in, headheight=14pt]{geometry}
\usepackage{fontspec}
\setmainfont{DejaVu Serif}
\setsansfont{DejaVu Sans}
\setmonofont{DejaVu Sans Mono}

\usepackage{xcolor}
\usepackage{titlesec}
\usepackage{fancyhdr}
\usepackage{enumitem}
\usepackage{hyperref}
\usepackage{booktabs}
\usepackage{tabularx}
\usepackage{longtable}
\usepackage{colortbl}
\usepackage{multirow}
\usepackage{array}
\usepackage{parskip}
\usepackage{tcolorbox}
\usepackage{graphicx}
\usepackage{lastpage}

% ── Color Scheme: Medical / Health ─────────────────────────────────────────
\definecolor{primary}{HTML}{1565C0}       % Clinical blue
\definecolor{secondary}{HTML}{2E7D32}     % Life green
\definecolor{tertiary}{HTML}{455A64}      % Warm slate
\definecolor{accent}{HTML}{0277BD}        % Bright blue
\definecolor{lightprimary}{HTML}{E3F2FD}
\definecolor{lightsecondary}{HTML}{E8F5E9}
\definecolor{lighttertiary}{HTML}{ECEFF1}
\definecolor{rowgray}{HTML}{F5F5F5}
\definecolor{bordergray}{HTML}{BDBDBD}
\definecolor{subtlegray}{HTML}{757575}
\definecolor{darktext}{HTML}{212121}

% ── Section Formatting ──────────────────────────────────────────────────────
\titleformat{\section}
  {\Large\bfseries\sffamily\color{primary}}
  {\thesection}{1em}{}
  [\vspace{-0.5em}\textcolor{bordergray}{\rule{\textwidth}{0.4pt}}]

\titleformat{\subsection}
  {\large\bfseries\sffamily\color{secondary}}
  {\thesubsection}{1em}{}

\titleformat{\subsubsection}
  {\normalsize\bfseries\sffamily\color{tertiary}}
  {\thesubsubsection}{1em}{}

\titlespacing*{\section}{0pt}{1.5em}{0.8em}
\titlespacing*{\subsection}{0pt}{1.2em}{0.5em}
\titlespacing*{\subsubsection}{0pt}{0.8em}{0.3em}

% ── Custom Environments ─────────────────────────────────────────────────────
\newcommand{\stageheader}[2]{%
  \noindent\colorbox{#2}{\makebox[\textwidth][l]{%
    \textcolor{white}{\bfseries\sffamily\hspace{6pt}#1\hspace{6pt}}}}%
  \vspace{0.5em}\par%
}

\newtcolorbox{asciibox}{
  colback=rowgray, colframe=bordergray,
  arc=1pt, boxrule=0.5pt,
  left=6pt, right=6pt, top=4pt, bottom=4pt,
  fontupper=\small\ttfamily,
  breakable
}

\newtcolorbox{quotebox}{
  colback=lightprimary, colframe=primary,
  arc=2pt, boxrule=0.5pt,
  left=12pt, right=12pt, top=8pt, bottom=8pt,
  breakable
}

\newcommand{\metafield}[2]{\textbf{\sffamily #1} #2\par}

% ── Table Helpers ───────────────────────────────────────────────────────────
\newcolumntype{L}[1]{>{\raggedright\arraybackslash}p{#1}}
\newcolumntype{C}[1]{>{\centering\arraybackslash}p{#1}}
\newcommand{\tableheader}[1]{\textbf{\sffamily\textcolor{white}{#1}}}

% ── Headers / Footers ───────────────────────────────────────────────────────
\pagestyle{fancy}
\fancyhf{}
\fancyhead[L]{\small\sffamily\textcolor{subtlegray}{Personal Care: Physical \& Mental}}
\fancyhead[R]{\small\sffamily\textcolor{subtlegray}{GSD Mission Package}}
\fancyfoot[C]{\small\sffamily\textcolor{subtlegray}{Page \thepage\ of \pageref{LastPage}}}
\renewcommand{\headrulewidth}{0.4pt}
\renewcommand{\footrulewidth}{0pt}

% ── Hyperlinks ──────────────────────────────────────────────────────────────
\hypersetup{
  colorlinks=true,
  linkcolor=secondary,
  urlcolor=secondary,
  pdfauthor={GSD Ecosystem / Miles Tiberius Foxglove},
  pdftitle={Personal Care: Physical and Mental -- Mission Package},
  pdfsubject={Vision to Mission Pipeline},
}

% ── Document ────────────────────────────────────────────────────────────────
\begin{document}

%% ══════════════════════════════════════════════════════════════════
%%  TITLE PAGE
%% ══════════════════════════════════════════════════════════════════
\thispagestyle{empty}
\begin{center}
\vspace*{3cm}

{\small\sffamily\textcolor{primary}{\textbf{GSD RESEARCH MISSION PACKAGE}}}

\vspace{0.3cm}
\textcolor{bordergray}{\rule{8cm}{0.4pt}}

\vspace{1.5cm}
{\fontsize{28}{34}\selectfont\bfseries\sffamily\textcolor{primary}{Personal Care:\\[0.3em]Physical \& Mental}}

\vspace{0.8cm}
{\Large\sffamily\textcolor{secondary}{The Whole-System Human: Evidence-Based Self-Care}}

\vspace{0.5cm}
{\large\sffamily\itshape\textcolor{tertiary}{A Deep Research Study}}

\vspace{3cm}
{\sffamily\textcolor{accent}{Vision $\rightarrow$ Research $\rightarrow$ Mission}}\\[0.3em]
{\small\sffamily\textcolor{accent}{Full Pipeline Execution}}

\vspace{3cm}
{\small\sffamily\textcolor{subtlegray}{Date: April 5, 2026  |  Status: Ready for GSD Orchestrator}}\\[0.3em]
{\small\sffamily\textcolor{subtlegray}{Activation Profile: Squadron  |  Estimated: 5--8 context windows across 3--4 sessions}}
\end{center}
\newpage

%% ══════════════════════════════════════════════════════════════════
%%  TABLE OF CONTENTS
%% ══════════════════════════════════════════════════════════════════
\tableofcontents
\newpage

%% ══════════════════════════════════════════════════════════════════
%%  STAGE 1 — VISION DOCUMENT
%% ══════════════════════════════════════════════════════════════════
\stageheader{STAGE 1 --- VISION DOCUMENT}{primary}

\section{Personal Care: Physical \& Mental --- Vision Guide}

\metafield{Date:}{April 5, 2026}
\metafield{Status:}{Initial Vision / Pre-Research Verification Complete}
\metafield{Depends On:}{GSD Ecosystem Core, Lifestyle Medicine Framework (ACLM/NIH)}
\metafield{Context:}{Cold-start deep research mission --- no prior conversation context harvested}

\subsection{Vision}

The engineering career teaches a specific truth: a system cannot sustain output beyond the quality of its power supply. You can optimize the pipeline, tune the chipset, and refine every instruction path --- but if the substrate is running hot, starved, or fragmented, the architecture collapses regardless. Personal care is the power supply for every other endeavor. It is not a luxury addendum to productivity. It is the floor on which everything else rests.

The GSD prime directive --- giving people their lives back --- begins here. Not with automation, not with agent orchestration, not with mission pipelines. It begins with the recognition that no amount of elegant tooling compensates for a human operating on three hours of sleep, a processed-food diet, zero physical movement, and social isolation. The Amiga Principle applies: specialized, high-quality input to the human system produces outsized output. Brute-force hours do not.

This mission investigates personal care across its five interlocking pillars: physical foundations (movement, nutrition, sleep), mental and emotional health, social connection, lifestyle systems and habits, and the mind-body interface that weaves them together. The research follows the evidence --- primarily NIH, CDC, WHO, and peer-reviewed clinical science --- not wellness marketing. The goal is a practical, actionable knowledge base that a builder, a parent, or anyone trying to reclaim their time and health can actually use.

The through-line is simple: the same architectural clarity that produces elegant software produces a well-functioning human system. Understand the dependencies. Reduce the friction. Give the system what it actually needs.

\subsection{Problem Statement}

The following five problems are well-documented, measurable, and directly addressable by this research:

\begin{enumerate}[leftmargin=1.5em]
  \item \textbf{The disconnection between knowledge and behavior.} Most people know they should sleep more, move more, and eat better. The gap between knowing and doing is a system design problem, not a willpower problem. Evidence-based habit architecture closes this gap.

  \item \textbf{Fragmentation of physical and mental health.} Clinical systems treat body and mind as separate departments. Yet depression increases cardiovascular risk; chronic pain elevates anxiety; poor sleep degrades emotional regulation. Treating pillars in isolation produces incomplete outcomes.

  \item \textbf{The loneliness epidemic as an unaddressed health variable.} The 2023 U.S. Surgeon General's Advisory and the 2025 WHO Commission on Social Connection both classify social isolation as a major public health threat --- comparable in mortality risk to smoking. Most personal care frameworks omit it entirely.

  \item \textbf{Self-care frameworks optimized for marketing rather than evidence.} The wellness industry generates billions in revenue from supplements, apps, and protocols with weak scientific support. A rigorous, source-anchored reference is needed to distinguish signal from noise.

  \item \textbf{Sustainability failure in personal care routines.} Most interventions produce short-term changes that do not persist. The science of behavior change, habit formation, and environmental design points to specific, replicable patterns for durable change --- patterns not widely taught.
\end{enumerate}

\subsection{Core Concept}

\textbf{One-line model:} The human body-mind is an interdependent system; targeted inputs to any one pillar propagate improvements across all others.

This is not holism as metaphor. It is measurable systems physiology. Exercise improves sleep quality, which improves emotional regulation, which improves social behavior, which reduces loneliness-driven cortisol, which improves immune function, which enables more exercise. The pillars are not additive --- they are multiplicative when properly aligned.

\subsubsection{Domain Architecture}

\begin{asciibox}
PERSONAL CARE SYSTEM ARCHITECTURE
==================================

        [ PHYSICAL FOUNDATIONS ]
        Movement | Nutrition | Sleep
               |          |
               v          v
    [ MIND-BODY INTERFACE ] <--------> [ MENTAL / EMOTIONAL HEALTH ]
    Gut-Brain Axis                      Stress | Cognition | Emotion
    Psychoneuroimmunology
    HPA Axis / Cortisol
               |
               v
    [ SOCIAL CONNECTION ]
    Community | Belonging | Purpose
               |
               v
    [ LIFESTYLE SYSTEMS ]
    Habits | Environment | Routines | Identity
               |
               v
    [ WHOLE-SYSTEM OUTPUT ]
    Sustained Energy | Emotional Resilience
    Cognitive Clarity | Creative Capacity
    Longevity | Presence for Others
\end{asciibox}

\subsubsection{Module Architecture}

\begin{asciibox}
MODULE STRUCTURE (5 Modules, Squadron Profile)
===============================================

  Module A: Physical Foundations
    - Movement / Exercise Science
    - Nutritional Science (macro + micro)
    - Sleep Architecture + Circadian Biology

  Module B: Mental & Emotional Health
    - Stress Physiology + Regulation
    - Cognitive Health + Neuroplasticity
    - Emotional Regulation Frameworks

  Module C: Social Connection
    - Loneliness Epidemiology
    - Social Health as Clinical Variable
    - Community and Belonging Interventions

  Module D: Lifestyle Systems & Habit Architecture
    - Behavior Change Science
    - Environmental Design
    - Identity-Based Habit Formation

  Module E: Mind-Body Interface
    - Gut-Brain Axis + Microbiome
    - Psychoneuroimmunology
    - HPA Axis / Allostatic Load

  Cross-Module Synthesis: Integration Report
    - Pillar Interaction Model
    - Practical Protocol Synthesis
    - Personalization Framework

  Dependency Flow:
    A + B (parallel) -> E -> C -> D -> Synthesis
\end{asciibox}

\subsection{Research Modules}

\subsubsection{Module A: Physical Foundations}
Surveys the three primary physical pillars: exercise and movement science (dose-response relationships, modality comparisons, minimal effective dose); nutritional science anchored in Dietary Guidelines for Americans 2020--2025, Mediterranean and whole-food approaches, micronutrient sufficiency; and sleep architecture including circadian biology, sleep stages, NIH/NHLBI adult sleep duration guidelines (7--9 hours), and the bidirectional relationship between sleep and all other health variables.

\subsubsection{Module B: Mental and Emotional Health}
Covers stress physiology (HPA axis, cortisol dynamics, allostatic load theory), cognitive health and neuroplasticity (evidence-based practices that preserve and enhance cognition across the lifespan), and emotional regulation frameworks including evidence from clinical psychology and the NIH's Research Domain Criteria (RDoC).

\subsubsection{Module C: Social Connection}
Investigates the scientific evidence on social isolation and loneliness as health determinants, drawing on the 2024 World Psychiatry meta-review (Holt-Lunstad), the 2025 WHO Commission on Social Connection flagship report, and the U.S. Surgeon General's 2023 Advisory. Identifies intervention approaches at individual, community, and systems levels.

\subsubsection{Module D: Lifestyle Systems and Habit Architecture}
Synthesizes the science of behavior change --- the six pillars of lifestyle medicine (nutrition, exercise, relationships, stress, sleep, substance use), habit formation research, environmental design principles, and identity-based change. Sources include the American College of Lifestyle Medicine (ACLM), NIH behavioral intervention literature, and clinical trial evidence.

\subsubsection{Module E: Mind-Body Interface}
Examines the mechanistic pathways connecting mental and physical states: gut-brain axis and microbiome-mood connections, psychoneuroimmunology (how psychological states alter immune function), and HPA axis regulation. This module provides the biological substrate explaining why the other pillars interact.

\subsection{Chipset Configuration}

\begin{asciibox}
CHIPSET CONFIGURATION (Amiga Metaphor)
=======================================

  Agnus   [Opus]   -- Synthesis + Cross-Module Integration
                      Judgment-heavy: pillar interaction model,
                      personalization framework, through-line
                      narrative. ~30% of context budget.

  Denise  [Sonnet] -- Research Compilation + Document Assembly
                      Module surveys A-E, source verification,
                      structured content, wave-by-wave output.
                      ~60% of context budget.

  Paula   [Haiku]  -- Scaffolding + Shared Schemas
                      Shared source index, templates, test
                      harnesses, bibliography formatting.
                      ~10% of context budget.

  Model Split: 30% Opus / 60% Sonnet / 10% Haiku
\end{asciibox}

\subsection{Scope Boundaries}

\subsubsection{In Scope (v1.0)}
\begin{itemize}[leftmargin=1.5em]
  \item Evidence-based personal care across all five pillars for generally healthy adults
  \item NIH, CDC, WHO, and peer-reviewed journal sources only
  \item Practical, actionable synthesis alongside research findings
  \item Interaction effects between pillars (cross-module integration)
  \item Habit formation and behavior change science
  \item Social connection as a clinical health variable
  \item Mind-body interface mechanisms (gut-brain, psychoneuroimmunology)
\end{itemize}

\subsubsection{Out of Scope}
\begin{itemize}[leftmargin=1.5em]
  \item Clinical treatment protocols for diagnosed mental illness (refer to licensed professionals)
  \item Supplement or pharmaceutical recommendations
  \item Age-specific pediatric or geriatric clinical protocols (general principles only)
  \item Wellness marketing content or unverified interventions
  \item Genetic or precision-medicine personalization (v2.0 candidate)
\end{itemize}

\subsection{Success Criteria}

\begin{enumerate}[leftmargin=1.5em]
  \item All factual claims attributed to NIH, CDC, WHO, peer-reviewed journals, or ACLM sources.
  \item Physical Foundations module covers all three pillars (exercise, nutrition, sleep) with dose-response data.
  \item Sleep module includes NHLBI adult duration guidelines and circadian biology fundamentals.
  \item Social connection module integrates the 2025 WHO Commission findings and 2023 Surgeon General Advisory.
  \item Mind-Body Interface module identifies at least three distinct mechanistic pathways.
  \item Habit Architecture module includes at least one evidence-based behavior change model (e.g., lifestyle medicine six pillars, implementation intentions, identity-based habits).
  \item Cross-module synthesis produces a practical pillar interaction model.
  \item All five modules are independently verifiable (each cites primary sources).
  \item Safety-critical tests pass: no unsourced clinical claims; no contraindicated advice for general audiences.
  \item Personalization framework provides differentiation criteria without requiring clinical assessment.
  \item Final document is suitable for direct use by a non-specialist reader without additional context.
  \item LaTeX compiles cleanly in three XeLaTeX passes; PDF renders without errors.
\end{enumerate}

\subsection{Relationship to Other GSD Documents}

\begin{center}
\rowcolors{2}{rowgray}{white}
\begin{tabularx}{\textwidth}{L{5cm}L{4cm}X}
\toprule
\rowcolor{primary}
\tableheader{Document} & \tableheader{Relationship} & \tableheader{Notes} \\
\midrule
The Space Between (Foxglove) & Philosophical parent & Mathematical and philosophical spine; ``spaces between'' as navigational metaphor applies directly to the pauses that enable recovery \\
Fox Infrastructure Group Vision & Organizational context & A 48--100 year Pacific Rim infrastructure org requires sustained human performance; personal care is an organizational asset \\
GSD Skill Creator v1.49 & Execution environment & Mission package handed to skill-creator for orchestrated execution \\
College of Knowledge (v1.50 target) & Curriculum sibling & Personal care science is a foundational department in the College of Knowledge \\
GSD Ecosystem Field Guide & Ecosystem reference & Amiga Principle, chipset model, and activation profiles formalized there \\
\bottomrule
\end{tabularx}
\end{center}

\subsection{The Through-Line}

In electronics, the power supply rail is the quietest, least glamorous component of any system. It has no user interface. It produces no visible output. But degrade it, and every downstream component performs below specification. Overclock the CPU on a failing PSU and you get erratic behavior, thermal runaway, eventual failure.

The human body-mind is no different. The work of personal care is the work of maintaining the rail. Sleep is not laziness; it is active neural consolidation and physical repair. Movement is not vanity; it is the signal that tells the body it is still needed and vital. Nutrition is not restriction; it is choosing the substrate that the system was designed to run on. Social connection is not soft; it is, by the evidence, as mortality-significant as smoking. And the mind-body interface that weaves them together is not mysticism --- it is measurable physiology.

The GSD vision --- giving people their lives back --- cannot be achieved by tools alone. It must also be achieved by helping people occupy their own lives with full capacity. This mission is that work: rigorous, sourced, actionable, and built to last.

\newpage

%% ══════════════════════════════════════════════════════════════════
%%  STAGE 2 — RESEARCH REFERENCE
%% ══════════════════════════════════════════════════════════════════
\stageheader{STAGE 2 --- RESEARCH REFERENCE}{secondary}

\section{Personal Care: Physical \& Mental --- Research Reference}

\subsection{How to Use This Document}

This reference compiles findings from authoritative sources organized by module. Each subsection is self-contained and citable. Agents assigned to a specific module should read the corresponding section in full before beginning output generation.

\textbf{Key source organizations:}
\begin{itemize}[leftmargin=1.5em]
  \item National Institutes of Health (NIH) --- nih.gov
  \item Centers for Disease Control and Prevention (CDC) --- cdc.gov
  \item National Heart, Lung, and Blood Institute (NHLBI) --- nhlbi.nih.gov
  \item World Health Organization (WHO) --- who.int
  \item American College of Lifestyle Medicine (ACLM) --- lifestylemedicine.org
  \item National Institute of Mental Health (NIMH) --- nimh.nih.gov
  \item U.S. Surgeon General's Office --- hhs.gov/surgeongeneral
  \item PubMed / PMC (National Library of Medicine) --- pubmed.ncbi.nlm.nih.gov
\end{itemize}

\subsection{Module A Reference: Physical Foundations}

\subsubsection{Movement and Exercise Science}

Physical inactivity is one of the leading modifiable risk factors for chronic disease globally. The CDC reports that only 50\% of U.S. adults meet recommended physical activity levels. Regular exercise produces documented benefits across cardiovascular, metabolic, musculoskeletal, cognitive, and mental health domains.

\textbf{Adult Physical Activity Guidelines (U.S. DHHS):}
\begin{itemize}[leftmargin=1.5em]
  \item 150--300 minutes per week of moderate-intensity aerobic activity, OR
  \item 75--150 minutes per week of vigorous-intensity aerobic activity
  \item Muscle-strengthening activities on 2 or more days per week
  \item Reducing sedentary time provides additional benefit independent of total activity
\end{itemize}

\textbf{Mental health effects of exercise (evidence-based):}
Physical activity has established effects on depression, anxiety, and cognitive function. Exercise is classified as a non-pharmacological treatment for multiple mental health conditions. Structured exercise programs have been shown to regulate mood, reduce anxiety symptoms, improve sleep quality, and enhance executive cognitive function.

\begin{center}
\rowcolors{2}{rowgray}{white}
\begin{tabularx}{\textwidth}{L{3.5cm}L{3.5cm}X}
\toprule
\rowcolor{secondary}
\tableheader{Exercise Type} & \tableheader{Primary Benefit} & \tableheader{Mental Health Linkage} \\
\midrule
Aerobic (walking, cycling) & Cardiovascular health & Depression, anxiety reduction \\
Resistance training & Musculoskeletal strength & Self-efficacy, depression \\
Flexibility / mobility & Joint health, injury prevention & Stress reduction, body awareness \\
Mind-body (yoga, tai chi) & Balance, proprioception & Anxiety, cortisol regulation \\
High-intensity interval (HIIT) & Metabolic efficiency & BDNF production, neuroplasticity \\
\bottomrule
\end{tabularx}
\end{center}

\subsubsection{Nutritional Science}

The CDC reports that only 10\% of U.S. adults consume the daily recommended amounts of fruits and vegetables, and that 90\% have diets high in sodium. Poor nutrition is a direct risk factor for diabetes, cardiovascular disease, stroke, obesity, and certain cancers. Emerging research also links dietary patterns to mental health outcomes, with several studies suggesting that Mediterranean-style diets reduce risk of depression and anxiety.

\textbf{Six evidence-based nutritional principles (NIH / USDA):}
\begin{enumerate}[leftmargin=1.5em]
  \item Emphasize fruits, vegetables, whole grains, lean proteins, and healthy fats
  \item Limit sodium, added sugars, saturated fats, and ultra-processed foods
  \item Maintain adequate micronutrient intake (particularly vitamin D, magnesium, B vitamins)
  \item Align eating timing with circadian rhythms (avoid late large meals)
  \item Prioritize dietary patterns over individual nutrients
  \item Sustainable dietary change outperforms short-term restriction in long-term outcomes
\end{enumerate}

\textbf{Diet-sleep-exercise interdependence:} Sleep Foundation research documents that diet and exercise are both independent predictors of sleep quality, and that sleep quality in turn influences food choice behavior and exercise motivation --- forming a triadic feedback system. Micronutrient inadequacy has been correlated with short sleep duration (NHANES 2005--2016 analysis, NIH).

\subsubsection{Sleep Architecture and Circadian Biology}

Sleep is a biologically active process critical to neurological, metabolic, immune, and endocrine function. The NHLBI defines adequate adult sleep as 7--9 hours per night. The CDC classifies insufficient sleep as a public health epidemic.

\textbf{Consequences of inadequate sleep (NHLBI):}
\begin{itemize}[leftmargin=1.5em]
  \item Increased risk of obesity, stroke, hypertension, and coronary artery disease
  \item Impaired glucose tolerance and altered hormonal levels
  \item Heightened emotional reactivity and reduced executive cognitive function
  \item Increased systemic inflammation
  \item Impaired memory consolidation and reduced attention capacity
\end{itemize}

\textbf{Sleep hygiene evidence base (NIH / Sleep Foundation):}
Consistent sleep and wake times regulate circadian rhythm. Caffeine consumption --- even in moderate amounts --- disrupts sleep onset and quality across multiple population studies (children, adolescents, adults, military personnel). Late large meals delay sleep onset. A dark, quiet, cool sleep environment promotes sleep continuity. These effects are reversible: evidence indicates that attaining adequate sleep reverses inflammation and cognitive impairment effects.

\textbf{Exercise-sleep bidirectional relationship:} Regular physical activity is associated with improved sleep efficiency and reduced sleep onset latency. Aerobic and resistance training are both associated with improved sleep duration and quality (Dolezal et al., 2017, Advances in Preventive Medicine).

\subsection{Module B Reference: Mental and Emotional Health}

\subsubsection{Stress Physiology}

Stress activates the hypothalamic-pituitary-adrenal (HPA) axis, triggering cortisol release. Acute stress is adaptive. Chronic stress produces allostatic load --- cumulative physiological wear from sustained activation --- increasing risk for cardiovascular disease, immune dysregulation, metabolic disorders, and psychiatric conditions. Effective stress management is classified as a core lifestyle medicine pillar by the ACLM.

\textbf{Evidence-based stress management strategies (CDC):}
\begin{itemize}[leftmargin=1.5em]
  \item Regular physical activity (direct cortisol regulation)
  \item Adequate, consistent sleep (HPA axis recovery)
  \item Social connection and support (cortisol buffering via oxytocin)
  \item Mindfulness and relaxation practices (parasympathetic activation)
  \item Time outdoors / nature exposure (validated stress reduction)
  \item Reducing sustained news/media consumption
\end{itemize}

\subsubsection{Cognitive Health and Neuroplasticity}

The brain retains neuroplastic capacity across the lifespan. Exercise, learning, sleep, and nutrition all influence the production of brain-derived neurotrophic factor (BDNF), which supports neuronal health and cognitive function. Major depressive disorder is the leading cause of disability worldwide (WHO/NCBI StatPearls). Cognitive health is bidirectionally linked to physical activity, sleep quality, and social engagement.

A 2024 Yale Journal of Biology and Medicine review confirmed that patient mindset significantly influences disease progression and recovery outcomes --- establishing psychological wellbeing as a clinical variable, not merely a quality-of-life adjunct.

\subsubsection{Emotional Regulation}

Emotional regulation capacity is linked to sleep quality, stress load, and social connection. Sleep deprivation produces heightened emotional reactivity measurably similar to clinical anxiety states. The lifestyle medicine framework identifies emotional wellness as the integrating variable across all other pillars: each pillar either supports or degrades the capacity to regulate emotional states.

\subsection{Module C Reference: Social Connection}

\subsubsection{Loneliness as a Health Variable}

Social connection is a fundamental biological need. The 2024 World Psychiatry meta-review (Holt-Lunstad) documents social connection as an independent predictor of mental and physical health, with some of the strongest evidence on all-cause mortality. The 2025 WHO Commission on Social Connection flagship report calls for treating social health with the same urgency as physical and mental health.

\textbf{Mortality-level findings:}
\begin{itemize}[leftmargin=1.5em]
  \item Social isolation and loneliness are associated with mortality risk comparable to smoking 15 cigarettes per day
  \item Loneliness is associated with a 29\% increased risk of coronary heart disease
  \item Lonely individuals experience amplified and prolonged stress responses (Kang et al., 2024)
  \item Social disconnection is associated with chronic pain, insomnia, and reduced immune function
\end{itemize}

\textbf{Epidemic-level trends:}

\begin{center}
\rowcolors{2}{rowgray}{white}
\begin{tabularx}{\textwidth}{L{4cm}X}
\toprule
\rowcolor{secondary}
\tableheader{Country / Body} & \tableheader{Policy Response to Loneliness Epidemic} \\
\midrule
United Kingdom & Minister of Loneliness appointed 2018; national strategy enacted \\
Japan & Minister of Loneliness appointed 2021 \\
South Korea & Monthly stipends for socially isolated youth introduced 2023 \\
United States & Surgeon General Advisory and national strategy framework 2023 \\
WHO & Commission on Social Connection; flagship report published June 2025 \\
\bottomrule
\end{tabularx}
\end{center}

\subsubsection{Intervention Evidence}

Group-based interventions consistently outperform individual-based interventions in reducing loneliness (digital interventions meta-analysis, ScienceDirect 2025). Psychological interventions targeting negative social cognitions and emotional skills show effectiveness for reducing loneliness at the individual level. Physical activity within social contexts addresses both movement and connection simultaneously, making it a high-leverage intervention.

\subsection{Module D Reference: Lifestyle Systems and Habit Architecture}

\subsubsection{The Six Pillars of Lifestyle Medicine}

The American College of Lifestyle Medicine defines six interdependent pillars of lifestyle health: nutrition, physical activity, sleep, stress management, social connection, and avoidance of risky substances. These are not independent interventions --- they form an integrated system where deficiency in any one pillar creates drag on all others.

NCBI StatPearls (Lifestyle Mental Wellbeing for the Primary Care Visit) documents that major depressive disorder has become the leading cause of disability worldwide, and that lifestyle medicine pillars are both risk factors and protective factors for mental illness, chronic disease, and quality of life.

\subsubsection{Behavior Change Science}

\textbf{Key principles from clinical behavioral research:}
\begin{itemize}[leftmargin=1.5em]
  \item \textbf{Implementation intentions:} Specifying when, where, and how to perform a behavior more than doubles follow-through rate compared to goal-setting alone
  \item \textbf{Environmental design:} Reducing friction for healthy behaviors and increasing friction for unhealthy ones produces behavior change without relying on willpower
  \item \textbf{Identity-based habits:} Behaviors congruent with self-identity are more durable than behaviors motivated by outcome goals alone
  \item \textbf{Start small / minimum effective dose:} Sustainable change requires entry points small enough that motivation is not the limiting factor
  \item \textbf{Social accountability:} Community and peer structures substantially increase adherence to lifestyle change programs
\end{itemize}

\subsubsection{The Four Pillars of Self-Care (Northwestern Medicine Framework)}

Northwestern Medicine's clinical self-care framework identifies four foundational pillars for daily practice: sleep, nutrition, physical activity, and leisure/recreation/recharge. The framework emphasizes that these are not secondary to other demands --- they are the precondition for meeting other demands effectively. Setting firm boundaries to protect these pillars and communicating them to social networks increases both adherence and relationship support.

\subsection{Module E Reference: Mind-Body Interface}

\subsubsection{Gut-Brain Axis}

The enteric nervous system (``second brain'') communicates bidirectionally with the central nervous system via the vagus nerve. Gut microbiome composition influences neurotransmitter production (including serotonin precursors), inflammatory cytokine levels, and stress response modulation. Research published in Cells (2024) identifies microbiome, metabolomics, hormones, and stress as interconnected mediators of mental health disorders.

Dietary patterns directly influence microbiome composition, creating a diet-gut-brain pathway where nutritional choices have measurable neurological and psychological effects extending beyond macronutrient and caloric considerations.

\subsubsection{Psychoneuroimmunology}

Psychological states measurably alter immune function. Positive emotional states enhance immune system activity (Dillon et al., Int J Psychiatry Med). Chronic stress and negative emotional states are associated with increased inflammatory markers, reduced natural killer cell activity, and elevated susceptibility to infection and autoimmune dysregulation. Conversely, interventions that reduce psychological distress (exercise, social connection, mindfulness) produce quantifiable immune improvements.

\subsubsection{HPA Axis and Allostatic Load}

The HPA axis stress response is adaptive in acute doses and damaging in chronic activation. Allostatic load --- the cumulative biological cost of chronic stress adaptation --- increases risk of cardiovascular disease, metabolic syndrome, immune dysfunction, and neurological degeneration. Sleep, exercise, social support, and mindfulness each reduce allostatic load through distinct but converging mechanisms.

\subsection{Safety and Sensitivity Considerations}

\begin{center}
\rowcolors{2}{rowgray}{white}
\begin{tabularx}{\textwidth}{L{4.5cm}L{2.5cm}X}
\toprule
\rowcolor{primary}
\tableheader{Area} & \tableheader{Boundary Type} & \tableheader{Guidance} \\
\midrule
Clinical mental illness & ABSOLUTE GATE & Never substitute for licensed professional assessment or treatment \\
Medication interactions & ABSOLUTE GATE & No pharmaceutical or supplement dosing recommendations \\
Clinical eating disorders & ABSOLUTE GATE & Nutritional content must not include restrictive framing; refer to NEDA/ANAD resources \\
Chronic disease management & GATE & General principles only; defer specific clinical protocols to care team \\
Substance use guidance & ANNOTATE & Present evidence without advocacy; note SAMHSA resources \\
Specific health metrics & GATE & Population-level ranges only; individual assessment requires clinical evaluation \\
\bottomrule
\end{tabularx}
\end{center}

\subsection{Source Bibliography}

\subsubsection{Government and Agency Sources}
\begin{itemize}[leftmargin=1.5em]
  \item CDC --- About Sleep. cdc.gov/sleep. Updated February 12, 2025.
  \item CDC --- Physical Activity Guidelines. cdc.gov/physicalactivity.
  \item NHLBI --- Sleep Deprivation and Deficiency. nhlbi.nih.gov/sleep-deprivation.
  \item NIH / NCBI --- Nursing Health Promotion: Healthy Lifestyle. NBK615357.
  \item NIH / NCBI StatPearls --- Lifestyle Mental Wellbeing for the Primary Care Visit. NBK587344.
  \item NIH / NCBI StatPearls --- Lifestyle Prevention Measures for the Clinic Visit. NBK587343.
  \item USDA --- MyPlate Nutritional Guidelines. myplate.gov. 2020--2025.
  \item U.S. DHHS --- Physical Activity Guidelines for Americans, 2nd Edition. 2018.
  \item U.S. Surgeon General --- Advisory on the Healing Effects of Social Connection and Community. 2023.
  \item WHO --- New Framework for Mental Health, Brain Health and Substance Use. October 2024.
  \item WHO Commission on Social Connection --- Flagship Report. June 30, 2025.
\end{itemize}

\subsubsection{Peer-Reviewed Research}
\begin{itemize}[leftmargin=1.5em]
  \item Holt-Lunstad, J. (2024). Social connection as a critical factor for mental and physical health. \textit{World Psychiatry}, 23: 312--332.
  \item Prabakar et al. (2024). The Power of Thought: Psychological Attentiveness in Patient Trajectories. \textit{Yale J Biol Med}, 97(3): 335--347.
  \item Dolezal et al. (2017). Interrelationship between Sleep and Exercise: A Systematic Review. \textit{Advances in Preventive Medicine}.
  \item Verma et al. (2024). Gut-Brain Axis: Role of Microbiome, Metabolomics, Hormones, and Stress in Mental Health Disorders. \textit{Cells}, 13(17).
  \item Kang et al. (2024). Loneliness and the experiences of everyday stress and stressor-related emotion. \textit{Stress and Health}.
  \item Shams-White et al. (2023). Healthy Eating Index-2020. \textit{J Acad Nutr Diet}, 123(9): 1280--1288.
  \item Digital interventions for loneliness --- Meta-analysis. \textit{ScienceDirect / Computers in Human Behavior}, 2025.
\end{itemize}

\subsubsection{Professional Organizations}
\begin{itemize}[leftmargin=1.5em]
  \item American College of Lifestyle Medicine (ACLM) --- lifestylemedicine.org
  \item American Psychological Association (APA) --- apa.org
  \item Sleep Foundation --- sleepfoundation.org
  \item Northwestern Medicine --- nm.org (Four Pillars of Self-Care framework)
  \item National Alliance for Eating Disorders --- allianceforeatingdisorders.com
\end{itemize}

\newpage

%% ══════════════════════════════════════════════════════════════════
%%  STAGE 3 — MISSION PACKAGE
%% ══════════════════════════════════════════════════════════════════
\stageheader{STAGE 3 --- MISSION PACKAGE}{tertiary}

\section{Milestone Specification}

\subsection{Mission Objective}

Produce a comprehensive, evidence-based reference document on personal care (physical and mental) using authoritative sources exclusively. The output will be a structured research compendium across five interlocking modules, culminating in a cross-module synthesis and practical protocol guide, ready for publication at \texttt{tibsfox.com/Missions/} and direct use by the target reader.

\subsection{Architecture Overview}

\begin{asciibox}
WAVE EXECUTION ARCHITECTURE
============================

WAVE 0 [Sequential, Haiku]
  Foundation: Shared schemas, source index, test templates

WAVE 1A [Parallel, Sonnet]          WAVE 1B [Parallel, Sonnet]
  Module A: Physical Foundations      Module B: Mental / Emotional Health
  Module C: Social Connection         Module E: Mind-Body Interface

WAVE 2A [Sequential, Sonnet]
  Module D: Lifestyle Systems & Habit Architecture
  (depends on A, B, C, E outputs)

WAVE 3 [Sequential, Opus]
  Cross-Module Synthesis
  Pillar Interaction Model
  Practical Protocol Guide

WAVE 4 [Sequential, Sonnet]
  Publication Assembly
  LaTeX PDF Compilation
  Index Page Generation
  Verification Pass
\end{asciibox}

\subsection{System Layers}

\begin{enumerate}[leftmargin=1.5em]
  \item \textbf{Foundation Layer} --- Shared schemas, source verification index, test harness templates (Haiku)
  \item \textbf{Survey Layer} --- Five independent module research documents (Sonnet, parallel)
  \item \textbf{Synthesis Layer} --- Cross-module integration, interaction model, protocol guide (Opus)
  \item \textbf{Publication Layer} --- Final document assembly, LaTeX compilation, index page (Sonnet)
  \item \textbf{Verification Layer} --- Source audit, safety-critical test pass, criteria check (Sonnet)
\end{enumerate}

\subsection{Deliverables}

\begin{center}
\rowcolors{2}{rowgray}{white}
\begin{tabularx}{\textwidth}{C{0.5cm}L{4cm}L{4cm}X}
\toprule
\rowcolor{tertiary}
\tableheader{\#} & \tableheader{Deliverable} & \tableheader{Acceptance Criterion} & \tableheader{Spec} \\
\midrule
1 & Module A Survey & All three physical pillars covered with dose data & Sonnet, Wave 1A, 2000--3000 words \\
2 & Module B Survey & Stress, cognition, emotion with sourced findings & Sonnet, Wave 1A, 1500--2500 words \\
3 & Module C Survey & Loneliness epidemiology + interventions & Sonnet, Wave 1B, 1500--2000 words \\
4 & Module D Survey & Six pillars + behavior change science & Sonnet, Wave 2A, 2000--3000 words \\
5 & Module E Survey & Gut-brain, PNI, HPA axis mechanisms & Sonnet, Wave 1B, 1500--2000 words \\
6 & Synthesis Report & Pillar interaction model + practical protocol & Opus, Wave 3, 3000--4000 words \\
7 & LaTeX PDF & Compiles cleanly; all criteria verified & Sonnet, Wave 4 \\
8 & Index Page & HTML with download links and usage guide & Sonnet, Wave 4 \\
\bottomrule
\end{tabularx}
\end{center}

\subsection{Component Breakdown}

\begin{center}
\rowcolors{2}{rowgray}{white}
\begin{tabularx}{\textwidth}{L{3.5cm}L{3cm}C{1.5cm}C{2cm}}
\toprule
\rowcolor{tertiary}
\tableheader{Component} & \tableheader{Dependencies} & \tableheader{Model} & \tableheader{Est. Tokens} \\
\midrule
Source Index Schema & None & Haiku & 2,000 \\
Test Template Set & None & Haiku & 1,500 \\
Module A Survey & Source Index & Sonnet & 8,000 \\
Module B Survey & Source Index & Sonnet & 7,000 \\
Module C Survey & Source Index & Sonnet & 6,000 \\
Module E Survey & Source Index & Sonnet & 6,000 \\
Module D Survey & A, B, C, E & Sonnet & 8,000 \\
Cross-Module Synthesis & A, B, C, D, E & Opus & 12,000 \\
LaTeX Assembly & All Modules + Synthesis & Sonnet & 10,000 \\
Verification Pass & All outputs & Sonnet & 4,000 \\
Index Page & PDF + .tex & Sonnet & 2,000 \\
\midrule
\multicolumn{3}{r}{\textbf{Total Estimated:}} & \textbf{$\sim$66,500} \\
\bottomrule
\end{tabularx}
\end{center}

\subsection{Model Assignment Rationale}

\begin{center}
\rowcolors{2}{rowgray}{white}
\begin{tabularx}{\textwidth}{L{2cm}C{1.5cm}X}
\toprule
\rowcolor{tertiary}
\tableheader{Model} & \tableheader{Share} & \tableheader{Rationale} \\
\midrule
Opus & 30\% & Cross-module synthesis requires judgment: which findings are load-bearing vs. noise; how to frame the pillar interaction model for a non-specialist; how to write the practical protocol without overpromising or undersourcing \\
Sonnet & 60\% & Module surveys are structured compilation tasks: gather, organize, cite, format. High-quality output with clear schema. Verification pass is systematic, not judgmental \\
Haiku & 10\% & Schema and template generation is scaffolding work; no judgment required \\
\bottomrule
\end{tabularx}
\end{center}

\section{Wave Execution Plan}

\subsection{Wave Summary}

\begin{center}
\rowcolors{2}{rowgray}{white}
\begin{tabularx}{\textwidth}{C{1cm}L{3.5cm}C{2cm}C{2cm}X}
\toprule
\rowcolor{primary}
\tableheader{Wave} & \tableheader{Name} & \tableheader{Mode} & \tableheader{Model} & \tableheader{Output} \\
\midrule
0 & Foundation & Sequential & Haiku & Schemas, source index, test templates \\
1A & Physical + Mental & Parallel & Sonnet & Modules A, B \\
1B & Social + Interface & Parallel & Sonnet & Modules C, E \\
2 & Lifestyle Systems & Sequential & Sonnet & Module D \\
3 & Synthesis & Sequential & Opus & Integration report + protocol guide \\
4 & Publication & Sequential & Sonnet & PDF, HTML, verification \\
\bottomrule
\end{tabularx}
\end{center}

\subsection{Wave 0: Foundation}

\begin{center}
\rowcolors{2}{rowgray}{white}
\begin{tabularx}{\textwidth}{L{4cm}X}
\toprule
\rowcolor{tertiary}
\tableheader{Task} & \tableheader{Specification} \\
\midrule
Source Index Schema & JSON schema: source\_id, organization, url, pub\_date, claim\_types[], module\_applicability[] \\
Shared Safety Checklist & List of 12 safety-critical tests; format as YAML with test\_id, description, expected, severity \\
Module Output Template & Markdown template with: metadata header, findings sections, data tables, source citations \\
Test Harness & Checklist runner: verify each deliverable against success criteria \\
\bottomrule
\end{tabularx}
\end{center}

\subsection{Wave 1A: Physical Foundations and Mental Health (Parallel)}

\begin{center}
\rowcolors{2}{rowgray}{white}
\begin{tabularx}{\textwidth}{L{2cm}L{4cm}X}
\toprule
\rowcolor{primary}
\tableheader{Track} & \tableheader{Module} & \tableheader{Agent Instructions} \\
\midrule
Track A1 & Module A: Physical Foundations & Survey exercise guidelines (DHHS), nutrition (USDA/NIH), and sleep (NHLBI/CDC). Include dose-response tables. Cite all claims. Follow safety checklist. \\
Track A2 & Module B: Mental/Emotional & Survey stress physiology (HPA axis, allostatic load), cognitive health (BDNF, neuroplasticity), emotional regulation. Source: NIMH, Yale J Biol Med 2024, NCBI StatPearls. \\
\bottomrule
\end{tabularx}
\end{center}

\subsection{Wave 1B: Social Connection and Mind-Body Interface (Parallel)}

\begin{center}
\rowcolors{2}{rowgray}{white}
\begin{tabularx}{\textwidth}{L{2cm}L{4cm}X}
\toprule
\rowcolor{primary}
\tableheader{Track} & \tableheader{Module} & \tableheader{Agent Instructions} \\
\midrule
Track B1 & Module C: Social Connection & Survey loneliness epidemiology and mortality data. Cite Holt-Lunstad 2024, WHO Commission 2025, Surgeon General 2023. Document policy responses. Include intervention evidence. \\
Track B2 & Module E: Mind-Body Interface & Survey gut-brain axis (Verma 2024), psychoneuroimmunology (Dillon; Antoni), HPA axis / allostatic load mechanisms. Identify three or more distinct mechanistic pathways. \\
\bottomrule
\end{tabularx}
\end{center}

\subsection{Wave 2: Lifestyle Systems (Sequential)}

Agent reads all four completed Wave 1 module surveys, then produces Module D: Lifestyle Systems and Habit Architecture. Must reference ACLM six pillars framework, behavior change science (implementation intentions, environmental design, identity-based habits), and Northwestern Medicine's Four Pillars framework. Explicitly cross-references findings from Modules A, B, C, and E to show how systems interact.

\subsection{Wave 3: Synthesis (Sequential, Opus)}

Opus agent reads all five module surveys. Produces:
\begin{itemize}[leftmargin=1.5em]
  \item \textbf{Pillar Interaction Model:} Which pillar improvements cascade most strongly into others; which sequences are highest-leverage for a new practitioner
  \item \textbf{Practical Protocol Guide:} Minimum effective dose for each pillar; entry points for someone starting from zero; daily, weekly, and monthly cadences
  \item \textbf{Personalization Framework:} Non-clinical differentiation criteria (e.g., sedentary vs. active baseline; isolated vs. connected starting point; sleep-deficit vs. sleep-sufficient)
  \item \textbf{Gap and Uncertainty Summary:} Areas where evidence is robust vs. emerging vs. contested
\end{itemize}

\subsection{Wave 4: Publication and Verification (Sequential)}

\begin{itemize}[leftmargin=1.5em]
  \item Assemble all modules and synthesis into final LaTeX document
  \item Run three-pass XeLaTeX compilation
  \item Execute test harness against all 12 success criteria
  \item Generate index.html with download links
  \item Flag any test failures for human review (CAPCOM gate before delivery)
\end{itemize}

\subsection{Cache Optimization Strategy}

\begin{itemize}[leftmargin=1.5em]
  \item Source Index is compiled in Wave 0 and included at the top of every subsequent agent context --- maximizing cache hit rate across waves
  \item Module output templates are shared schemas --- agents produce consistent structure, reducing synthesis parsing overhead
  \item Wave 1A and 1B run in parallel, halving wall-clock time for survey compilation
  \item Each module survey is self-contained and cacheable; reruns of individual modules do not require re-running the full pipeline
  \item Opus synthesis context starts with a compressed summary of each module (500 tokens each) before loading full text --- preserving judgment capacity for integration work
\end{itemize}

\subsection{Token Budget}

\begin{center}
\rowcolors{2}{rowgray}{white}
\begin{tabularx}{\textwidth}{L{4cm}C{2cm}C{2cm}C{2cm}}
\toprule
\rowcolor{primary}
\tableheader{Stage} & \tableheader{Model} & \tableheader{Context Windows} & \tableheader{Est. Tokens} \\
\midrule
Wave 0: Foundation & Haiku & 1 & 3,500 \\
Wave 1A: Mod A + B & Sonnet & 2 & 15,000 \\
Wave 1B: Mod C + E & Sonnet & 2 & 12,000 \\
Wave 2: Mod D & Sonnet & 1 & 8,000 \\
Wave 3: Synthesis & Opus & 1--2 & 20,000 \\
Wave 4: Publication & Sonnet & 1 & 8,000 \\
\midrule
\textbf{Total} & & \textbf{8--9} & \textbf{$\sim$66,500} \\
\bottomrule
\end{tabularx}
\end{center}

\subsection{Dependency Graph}

\begin{asciibox}
DEPENDENCY GRAPH
=================

  [W0: Foundation]
        |
   _____|_____
  |           |
[W1A: A+B]  [W1B: C+E]    <-- parallel
  |     |    |     |
  |     +----+     |
  |         |      |
  +----+----+------+
       |
  [W2: Module D]            <-- sequential; depends on all W1
       |
  [W3: Synthesis]            <-- sequential; depends on D + all W1
       |
  [W4: Publication]          <-- sequential; terminal wave

  CAPCOM gates:
    Between W1 -> W2: human review of module quality
    Between W3 -> W4: human approval of synthesis framing
    W4 output: final human verification before delivery
\end{asciibox}

\section{Test Plan}

\subsection{Test Categories}

\begin{center}
\rowcolors{2}{rowgray}{white}
\begin{tabularx}{\textwidth}{L{4cm}C{2cm}C{2cm}X}
\toprule
\rowcolor{tertiary}
\tableheader{Category} & \tableheader{Count} & \tableheader{Priority} & \tableheader{Action on Fail} \\
\midrule
Safety-Critical & 4 & P0 (BLOCK) & Stop pipeline; require human review and correction \\
Core Functionality & 5 & P1 & Require correction before Wave 4 \\
Integration & 2 & P2 & Flag for synthesis agent; document in gap summary \\
Edge Cases & 2 & P3 & Document and annotate; do not block \\
\midrule
\textbf{Total} & \textbf{13} & & \\
\bottomrule
\end{tabularx}
\end{center}

\subsection{Safety-Critical Tests}

\begin{center}
\rowcolors{2}{rowgray}{white}
\begin{tabularx}{\textwidth}{C{1cm}L{4cm}X}
\toprule
\rowcolor{primary}
\tableheader{ID} & \tableheader{Verifies} & \tableheader{Expected Behavior} \\
\midrule
SC-01 & No unsourced clinical claims & Every specific factual claim cites NIH, CDC, WHO, ACLM, or peer-reviewed source by name \\
SC-02 & No clinical treatment recommendations & Document contains no specific dosing, prescription, or diagnosis-level guidance \\
SC-03 & No restrictive nutritional content & Nutritional section does not include caloric restriction framing, weight-loss targets, or content that could facilitate disordered eating \\
SC-04 & Mental illness referral gate & Any content touching clinical depression, anxiety disorder, or psychiatric conditions includes explicit ``consult a licensed professional'' language \\
\bottomrule
\end{tabularx}
\end{center}

\subsection{Verification Matrix}

\begin{center}
\rowcolors{2}{rowgray}{white}
\begin{tabularx}{\textwidth}{C{0.5cm}L{5cm}L{3cm}C{2cm}}
\toprule
\rowcolor{tertiary}
\tableheader{\#} & \tableheader{Success Criterion} & \tableheader{Test IDs} & \tableheader{Status} \\
\midrule
1 & All claims attributed to approved sources & SC-01 & Gate \\
2 & Physical module covers all three pillars & CORE-01 & Verify \\
3 & Sleep module includes NHLBI guidelines & CORE-02 & Verify \\
4 & Social module integrates 2025 WHO + 2023 SG & CORE-03 & Verify \\
5 & Mind-body module: 3+ mechanistic pathways & CORE-04 & Verify \\
6 & Habit module: 1+ evidence-based change model & CORE-05 & Verify \\
7 & Cross-module synthesis produced & INT-01 & Verify \\
8 & All five modules independently verifiable & INT-02 & Verify \\
9 & Safety-critical tests pass & SC-01--04 & BLOCK \\
10 & Personalization framework present & CORE-05 & Verify \\
11 & Reader-accessible without additional context & EDGE-01 & Review \\
12 & LaTeX compiles cleanly, 3 passes & EDGE-02 & Verify \\
\bottomrule
\end{tabularx}
\end{center}

\section{Mission Crew Manifest}

\begin{center}
\rowcolors{2}{rowgray}{white}
\begin{tabularx}{\textwidth}{L{3.5cm}C{1.5cm}X}
\toprule
\rowcolor{primary}
\tableheader{Role} & \tableheader{Agent} & \tableheader{Primary Responsibility} \\
\midrule
Mission Commander & Opus & Synthesis authority; final integration judgment \\
Research Lead A & Sonnet & Module A: Physical Foundations survey \\
Research Lead B & Sonnet & Module B: Mental / Emotional Health survey \\
Research Lead C & Sonnet & Module C: Social Connection survey \\
Research Lead D & Sonnet & Module D: Lifestyle Systems survey \\
Research Lead E & Sonnet & Module E: Mind-Body Interface survey \\
Integration Analyst & Opus & Cross-module pillar interaction model \\
Protocol Architect & Opus & Practical protocol guide and personalization framework \\
Document Assembler & Sonnet & LaTeX assembly, three-pass compilation \\
Verification Officer & Sonnet & Test harness execution, safety-critical gate \\
Foundation Engineer & Haiku & Schemas, source index, templates \\
Pub Engineer & Sonnet & Index page, tibsfox.com upload preparation \\
\bottomrule
\end{tabularx}
\end{center}

\section{Execution Summary}

\begin{center}
\rowcolors{2}{rowgray}{white}
\begin{tabularx}{\textwidth}{L{5cm}X}
\toprule
\rowcolor{secondary}
\tableheader{Metric} & \tableheader{Value} \\
\midrule
Activation Profile & Squadron (12 roles, 5 modules) \\
Total Modules & 5 + 1 synthesis \\
Deliverables & 8 (5 module surveys, synthesis, PDF, index.html) \\
Parallel Tracks & 2 (Wave 1A and 1B simultaneous) \\
CAPCOM Gates & 3 (Wave 1 review, Wave 3 approval, final delivery) \\
Safety-Critical Tests & 4 (all P0 / BLOCK) \\
Total Tests & 13 \\
Estimated Context Windows & 8--9 \\
Estimated Token Budget & 66,500 \\
Model Split & 30\% Opus / 60\% Sonnet / 10\% Haiku \\
Primary Sources & NIH, CDC, WHO, ACLM, NIMH, Surgeon General \\
Publication Target & tibsfox.com/Missions/ \\
LaTeX Toolchain & XeLaTeX + DejaVu fonts + 3-pass compilation \\
\bottomrule
\end{tabularx}
\end{center}

\vspace{2em}
\begin{quotebox}
\textit{``The power supply rail is the quietest, least glamorous component of any system. It has no user interface. It produces no visible output. But degrade it, and every downstream component performs below specification.''}

\vspace{0.5em}
{\small\sffamily\textcolor{subtlegray}{The through-line: personal care is the rail. Everything else --- the work, the relationships, the creative output, the mission --- runs on it. The GSD ecosystem gives people their time back. Personal care gives them the capacity to use it.}}
\end{quotebox}

\end{document}
